\documentclass[11pt]{article}

\usepackage[T1]{fontenc}
\usepackage[a4paper,margin=2.7cm]{geometry}
\usepackage{amsmath,amssymb,amsthm,mathtools}
\usepackage{xcolor}
\usepackage{hyperref}
\usepackage{microtype}

\hypersetup{
  colorlinks=true,
  linkcolor=blue!55!black,
  citecolor=blue!55!black,
  urlcolor=blue!55!black
}

\newtheorem{theorem}{Theorem}[section]
\newtheorem{lemma}[theorem]{Lemma}
\newtheorem{corollary}[theorem]{Corollary}
\theoremstyle{remark}
\newtheorem*{remark}{Remark}

\newcommand{\N}{\mathbb{N}}
\newcommand{\Q}{\mathbb{Q}}
\newcommand{\Z}{\mathbb{Z}}

\title{A Lean-verified offset construction\\
for the Erdős--Straus equation}
\author{}
\date{}

\begin{document}
\maketitle

\begin{abstract}
We give a self-contained formulation of an offset construction for the
Erdős--Straus equation
\[
  \frac4n=\frac1x+\frac1y+\frac1z.
\]
The construction packages the classical divisor split into a congruence
condition indexed by a prime \(p\equiv3\pmod4\).  Its full positive-integer
statement is formalised in Lean~4 with Mathlib: the development constructs the
denominators by natural-number division, proves the divisions exact and the
denominators positive, and concludes the equality in \(\Q\).  We also formalise
a fixed-divisor lifting theorem and three groups of explicit infinite families,
including the progressions arising from \(p=3\) and \(d_1=2,5\).  We make no
claim toward resolving the Erdős--Straus conjecture; the underlying divisor
split is classical.
\end{abstract}

\section{The divisor split}

The Erdős--Straus conjecture asks whether, for every integer \(n\ge2\), there
exist \(x,y,z\in\N_{>0}\) satisfying
\begin{equation}\label{eq:ES}
  \frac4n=\frac1x+\frac1y+\frac1z.
\end{equation}
The conjecture remains open.  We use the following classical algebraic
reduction; see, for example, Mordell~\cite{Mordell}.

Fix \(x\in\N_{>0}\) with \(4x>n\), and reduce
\[
  \frac4n-\frac1x=\frac{4x-n}{nx}=\frac ab,
  \qquad \gcd(a,b)=1.
\]
Multiplication and rearrangement give
\begin{equation}\label{eq:split}
  (ay-b)(az-b)=b^2.
\end{equation}

\begin{lemma}[positive divisor split]\label{lem:split}
Let \(n,x,a,b,d_1,d_2,y,z\in\N_{>0}\).  Suppose
\[
  a=4x-n,\qquad b=nx,\qquad d_1d_2=b^2,
\]
and
\[
  ay=b+d_1,\qquad az=b+d_2.
\]
Then \((x,y,z)\) satisfies \eqref{eq:ES}.
\end{lemma}

\begin{proof}
The last three equations are precisely the factorisation
\eqref{eq:split}.  Expanding it and using \(a=4x-n\) and \(b=nx\) gives
\[
  4xyz=n(xy+xz+yz).
\]
All four variables \(n,x,y,z\) are nonzero, so division by \(nxyz\) gives
\eqref{eq:ES}.
\end{proof}

The Lean development separates these two steps.  It first proves the
cross-multiplied identity over \(\Z\), then proves that a positive instance over
\(\N\) gives the rational equality in \(\Q\).

\section{A positive certificate}

The offset form avoids subtraction and is convenient both mathematically and
formally.

\begin{theorem}[positive offset certificate]\label{thm:certificate}
Let \(n,x,p,b,d_1,d_2,y,z\in\N\), with \(n,x,p,y,z>0\).  If
\begin{align}
  4x &= n+p,                 & b &= nx, \label{eq:certificate-a}\\
  d_1d_2 &= b^2,             & py &= b+d_1,\qquad pz=b+d_2, \label{eq:certificate-b}
\end{align}
then \((x,y,z)\) satisfies \eqref{eq:ES}.
\end{theorem}

\begin{proof}
The equations are the positive divisor split with \(a=p=4x-n\).  Equivalently,
substitution in the cross-multiplied target leaves the multiple
\[
  p\,(d_1d_2-b^2),
\]
which vanishes by \eqref{eq:certificate-b}.  Positivity then permits passage to
the rational identity.
\end{proof}

In Lean, \texttt{IsES n x y z} contains both the four positivity assertions and
the equality \eqref{eq:ES} in \(\Q\).  Theorem~\ref{thm:certificate} is the
declaration \texttt{isES\_of\_offset\_certificate}.

\section{The constructive offset theorem}

\begin{theorem}[constructive offset theorem]\label{thm:offset}
Let \(n,p,d_1\in\N\) satisfy
\[
  n\ge2,\qquad p\ \text{prime},\qquad
  p\equiv3\pmod4,\qquad n\equiv1\pmod4,\qquad p\nmid n.
\]
Set
\[
  x=\frac{n+p}{4},\qquad b=nx.
\]
Suppose \(d_1>0\), \(d_1\mid b\), and
\begin{equation}\label{eq:first-congruence}
  p\mid b+d_1,
\end{equation}
equivalently \(d_1\equiv-b\pmod p\).  Define
\begin{equation}\label{eq:witnesses}
  d_2=\frac{b^2}{d_1},\qquad
  y=\frac{b+d_1}{p},\qquad
  z=\frac{b+d_2}{p}.
\end{equation}
Then \(x,y,z\in\N_{>0}\), all divisions in \eqref{eq:witnesses} are exact, and
\((x,y,z)\) satisfies \eqref{eq:ES}.
\end{theorem}

\begin{proof}
The congruences on \(n\) and \(p\) show that \(4\mid n+p\), so \(x\in\N\) and
\(4x=n+p\).  Moreover \(x>0\).  If \(p\mid x\), then
\(p\mid4x=n+p\), hence \(p\mid n\), a contradiction.  Thus \(p\nmid b=nx\).

Since \(d_1\mid b\), the quotient \(d_2=b^2/d_1\) is exact and positive.  It
remains to prove \(p\mid b+d_2\).  Work in the field \(\Z/p\Z\).
Condition \eqref{eq:first-congruence} gives
\[
  d_1=-b.
\]
The element \(b\) is nonzero, hence so is \(d_1\).  From \(d_1d_2=b^2\), we may
cancel \(d_1\) to obtain
\[
  d_2=-b,
\]
and therefore \(p\mid b+d_2\).  Both definitions of \(y,z\) in
\eqref{eq:witnesses} are consequently exact.  Their numerators are positive,
so \(y,z>0\).  Theorem~\ref{thm:certificate} now applies.
\end{proof}

\begin{remark}
The primality assumption is used only for cancellation modulo \(p\).  The Lean
proof implements this step in \texttt{ZMod p}; the remainder of the construction
uses natural numbers, with the final equality stated over the rationals.
\end{remark}

\section{Fixed-divisor lifting}

For fixed \(p\), define
\[
  X_p(n)=\frac{n+p}{4},
  \qquad
  B_p(n)=nX_p(n),
\]
where division is natural-number division.  Call a triple \((p,d,n)\)
\emph{offset-admissible} if
\[
\begin{gathered}
  n\ge2,\quad p\ \text{prime},\quad
  p\equiv3\pmod4,\quad n\equiv1\pmod4,\quad p\nmid n,\\
  d>0,\quad d\mid B_p(n),\quad p\mid B_p(n)+d.
\end{gathered}
\]
These are exactly the hypotheses of Theorem~\ref{thm:offset}, packaged with a
fixed divisor \(d\).

\begin{theorem}[fixed-divisor lifting]\label{thm:lifting}
If \((p,d,n)\) is offset-admissible, then for every \(k\in\N\),
\[
  n_k=n+4pdk
\]
is offset-admissible for the same \(p,d\).  Consequently, on setting
\[
  x_k=X_p(n_k),\qquad b_k=B_p(n_k),\qquad
  d_{2,k}=\frac{b_k^2}{d},
\]
the positive natural numbers
\[
  y_k=\frac{b_k+d}{p},
  \qquad
  z_k=\frac{b_k+d_{2,k}}{p}
\]
satisfy
\[
  \frac4{n_k}=\frac1{x_k}+\frac1{y_k}+\frac1{z_k}.
\]
Moreover, \(k\mapsto n_k\) is strictly increasing, so these parameters form an
infinite arithmetic progression.
\end{theorem}

\begin{proof}
The congruences defining admissibility make \(n+p\) divisible by \(4\).
For \(n_k=n+4pdk\), exact division gives
\begin{equation}\label{eq:x-transport}
  X_p(n_k)=X_p(n)+pdk.
\end{equation}
Expanding \(B_p(n_k)=n_kX_p(n_k)\) using
\eqref{eq:x-transport} gives the exact identity
\begin{equation}\label{eq:b-transport}
  B_p(n_k)
  =
  B_p(n)+pdk\bigl(n+4X_p(n)+4pdk\bigr).
\end{equation}
The added term in \eqref{eq:b-transport} is divisible by both \(d\) and \(p\).
It follows respectively that \(d\mid B_p(n_k)\) and
\(p\mid B_p(n_k)+d\).  Also \(n_k\equiv n\pmod4\) and
\(n_k\equiv n\pmod p\), so the two modular hypotheses are preserved; all other
admissibility conditions are unchanged or immediate from \(n_k\ge n\).
Theorem~\ref{thm:offset} supplies the displayed positive representation,
including exactness of all divisions.  Finally,
\[
  n_{k+1}-n_k=4pd>0
\]
because a prime \(p\) and the admissible divisor \(d\) are positive.
\end{proof}

\section{Explicit infinite families}

We record the families formalised in the accompanying Lean file.  In every
corollary, \(k\in\N\).

\begin{corollary}[Family I]\label{cor:family-I}
Let
\[
  n=12k+9,\qquad x=3k+3,\qquad b=n(k+1).
\]
Then
\[
  \frac4n=\frac1x+\frac1{b+1}+\frac1{b(b+1)}.
\]
\end{corollary}

\begin{proof}
Here \(4x-n=3\) and \(\gcd(3,nx)=3\), so the reduced numerator is \(1\).
Taking \(d_1=1\) and \(d_2=b^2\) in the divisor split gives the displayed
denominators.  The Lean proof is a direct positive polynomial identity.
\end{proof}

\begin{corollary}[\(p=3,d_1=2\)]\label{cor:d2}
Let
\[
  n=24k+5\qquad\text{or}\qquad n=24k+13.
\]
Put \(x=(n+3)/4\) and \(b=nx\).  Then
\[
  \frac4n
  =\frac1x
   +\frac1{(b+2)/3}
   +\frac1{(b+b^2/2)/3}.
\]
\end{corollary}

\begin{proof}
The triples \((3,2,5)\) and \((3,2,13)\) are offset-admissible.  Applying
Theorem~\ref{thm:lifting} gives periods \(4\cdot3\cdot2=24\) and hence the two
displayed progressions.
\end{proof}

\begin{corollary}[\(p=3,d_1=5\)]\label{cor:d5}
Let
\[
  n\in
  \{60k+5,\ 60k+17,\ 60k+25,\ 60k+37\}.
\]
Put \(x=(n+3)/4\) and \(b=nx\).  Then
\[
  \frac4n
  =\frac1x
   +\frac1{(b+5)/3}
   +\frac1{(b+b^2/5)/3}.
\]
\end{corollary}

\begin{proof}
The triples with \(p=3,d_1=5\) and seeds
\[
  n\in\{5,17,25,37\}
\]
are offset-admissible.  Theorem~\ref{thm:lifting} gives period
\(4\cdot3\cdot5=60\), yielding the four displayed progressions.
\end{proof}

The two progressions in Corollary~\ref{cor:d2} are the nonmultiples of \(3\)
in the classical Mordell region \(n\equiv5\pmod8\); the remaining progression
\(n=24k+21\) is contained in Family~I.  The purpose of including them is to
exhibit the general construction and its formal specialisation, not to claim a
new residue family.

\section{Lean formalisation}

The formalisation is contained in \texttt{lean/ESTheorem.lean} and uses Lean~4
with Mathlib.  Its public positive interface consists of:
\begin{itemize}
  \item \texttt{IsES} and \texttt{isES\_of\_polynomial};
  \item \texttt{isES\_of\_offset\_certificate};
  \item \texttt{second\_divisibility};
  \item \texttt{constructive\_offset};
  \item the reducible definitions \texttt{offsetX}, \texttt{offsetB}, and the
        packaged interface \texttt{OffsetAdmissible} with
        \texttt{OffsetAdmissible.isES};
  \item \texttt{offsetX\_add\_period}, \texttt{offsetB\_add\_period}, and
        \texttt{OffsetAdmissible.add\_period};
  \item \texttt{fixed\_d\_progression} and
        \texttt{fixed\_d\_progression\_strictMono};
  \item \texttt{family\_I\_positive}, \texttt{family\_d2\_24k5},
        \texttt{family\_d2\_24k13}, and the four
        \texttt{family\_d5\_60k\(\cdot\)} declarations.
\end{itemize}
For every explicit progression, the file also checks the instances \(k=0\) and
\(k=1\).  The theorem statement includes the positive denominators and the
rational unit-fraction equality; the formalisation is therefore not merely a
check of a polynomial identity.

\section{Scope}

The divisor split is classical, and offset choices of \(x\) belong to the
standard toolkit for Egyptian-fraction problems.  The contribution of this note
is a compact theorem interface, an end-to-end Lean proof of the positive
construction, and a small collection of formally checked infinite
specialisations.  The full Erdős--Straus conjecture remains outside its scope.

\begin{thebibliography}{9}

\bibitem{Mordell}
L.~J.~Mordell,
\emph{Diophantine Equations},
Academic Press, 1969.

\bibitem{Vaughan}
R.~C.~Vaughan,
\emph{On a problem of Erdős, Straus and Schinzel},
J. London Math. Soc. (2) 2 (1970), 391--396.

\bibitem{mathlib}
The Mathlib Community,
\emph{Mathlib: the mathematical library of Lean~4},
\href{https://github.com/leanprover-community/mathlib4}{project repository}.

\end{thebibliography}

\end{document}
